\documentclass{article}

\usepackage[T1]{fontenc}

\usepackage{booktabs,graphicx,siunitx}

\begin{document}

\section{ANCOVA}

\begin{tabular}{llll}
\toprule
Group & Adj. M & SE & 95\% CI \\
\midrule
Blended & 65.37 & 0.71 & [63.97, 66.78] \\
Flipped & 70.95 & 0.71 & [69.54, 72.35] \\
Lecture & 60.38 & 0.71 & [58.98, 61.78] \\
\bottomrule
\end{tabular}

\begin{tabular}{lllllll}
\toprule
Source & SS & df & MS & F & p & partial $\eta$\textsuperscript{2} [95\% CI] \\
\midrule
pretest\_score & 2086.72 & 1 & 2086.72 & 51.49 & $<$.001 & 0.18 [0.11, 1.00] \\
teaching\_method & 4463.20 & 2 & 2231.60 & 55.07 & $<$.001 & 0.32 [0.24, 1.00] \\
\bottomrule
\end{tabular}

\begin{tabular}{lllll}
\toprule
Pair & Mdiff (adj.) & SE & padj & 95\% CI \\
\midrule
Blended - Flipped & $-$5.57 & 1.01 & $<$.001 & [$-$7.95, $-$3.20] \\
Blended - Lecture & 4.99 & 1.01 & $<$.001 & [2.62, 7.37] \\
Flipped - Lecture & 10.57 & 1.01 & $<$.001 & [8.19, 12.94] \\
\bottomrule
\end{tabular}

\includegraphics{figures/01_ancova/figure_adjusted-means.png}

Compares group means after statistically removing the influence of a numeric covariate. First confirm the factor$\times$covariate interaction is non-significant (homogeneity of slopes) - if it is significant, the adjusted-means interpretation is not valid and a different model is needed. Then read the Factor row's F/p for the adjusted group effect, and the adjusted means for the covariate-controlled group values. Your significance threshold ($\alpha$) is 0.05.

Controlling for the covariate, an ANCOVA gave F(2,236)=55.07, p $<$ .001, partial $\eta$\textsuperscript{2}=.32 [0.24, 1.00].

Statistical basis: ANCOVA (group means adjusted for a covariate) (Fisher, R. A. (1925). Statistical Methods for Research Workers. Oliver \& Boyd.); Type III sums of squares (car::Anova) (Fox J, Weisberg S (2019). An R Companion to Applied Regression, Third edition. Sage, Thousand Oaks CA.); Estimated marginal (covariate-adjusted) means (Lenth RV (2024). emmeans: Estimated Marginal Means, aka Least-Squares Means. R package.); Partial $\eta$\textsuperscript{2} effect-size benchmarks (Cohen, J. (1988). Statistical Power Analysis for the Behavioral Sciences (2nd ed.). Routledge.).

\end{document}